\section{Comparing the local factors at neighboring couplings}
\label{sec:local-parameter-comparison}
We now prove the compensated comparison used in the fixed-volume
continuity argument. The important simplification is to preserve the
physical fine field while changing the Gaussian coordinate. Thus the
incoming measurable activity is evaluated at the same point in the
two representations. Only its value difference, and not a derivative
of that activity, is needed.

Fix a derivative order $j\in\mathbb N$ and a fine torus. For a scalar
function of the coarse variation $y$, write
\[
 \|f\|_{j,y_0}=\max_{0\le r\le j}\|D^rf(y_0)\|.
\]
All derivatives below are taken on a common admissible local chart,
with its fixed normalized directions. In particular, the assertions
include existence and continuity of these derivatives near $y_0$.
We use the product estimate
\begin{equation}\label{eq:local-parameter-product-jet}
       \|fh\|_{j,y_0}\le 3^j\|f\|_{j,y_0}\|h\|_{j,y_0}.
\end{equation}
The Leibniz formula has $2^r$ terms at order $r$, so the slightly
larger factor $3^j$ is valid also for $j=0$.

The base representation supplies summable elementary prices
$u_i\ge0$, with $B=\sum_i u_i<\infty$, for its compensated factors
$w_i$. All logarithms here are real. Assume explicitly a common
base value bound $|\ell_i(y_0)|\le M$ for their real exponents;
in the construction this follows from the summable incoming logarithmic
prices and the held-field estimate. A bound on $e^{\ell_i}-1$ alone
would not control a large negative $\ell_i$. We also use the following local estimates already required by
the compensated construction: uniform jets of its enforcement
cutoffs; a bound $A s$ for the packet action jets on every packet of
size $s$ whose cutoff is not identically zero near $y_0$; and, for
its density compensation $K_I$ and quadratic exponent $q_I$, for
every subset $I$ of read bonds,
\begin{equation}\label{eq:local-parameter-density-input}
 \|K\|_{j,y_0}\le C_0F_0(z)^N,\qquad
 \|q\|_{j,y_0}\le C_0F_0(z)^2.
\end{equation}
Here $N\in\mathbb N$, $C_0\ge0$, and
$F_0(z)=1+d+\sum_{r=1}^d|z_r|\ge1$ uses the full finite Gaussian
coordinate space. The constants $N,C_0$ are common to every such $I$; we suppress that
index in the formulas. These estimates must hold for the same base
representation, forests and compensation maps. Their uniform
purchase in the scale construction is a separate requirement; the
comparison below uses their actual values at the fixed stage.

Let $R'$ be a neighboring representation with the same dictionaries,
paths and masks as $R$, and let $t>0$ bound its input differences as
in the fixed-volume continuity argument. Put
$\lambda=g(\beta)/g(\beta')$. The identity
$g(\beta')\lambda=g(\beta)$ implies that the fine field at $(R',\lambda z)$
is precisely the fine field at $(R,z)$. It also gives
$\beta\lambda^2=\beta'$ and makes the compensated noise course for
$R'$ equal to $\lambda$ times the course for $R$.

\subsection{Elementary exponential comparisons}
Two elementary facts explain the logarithmic and packet estimates.
First, suppose $\|e^f-1\|_{j,y_0}\le u$, $u\le B$, and
$|f(y_0)|\le M$. Then
\begin{equation}\label{eq:local-parameter-log-recovery}
                         \|f\|_{j,y_0}\le C_{j,M,B}u.
\end{equation}
Indeed $f=\log(1+(e^f-1))$. The derivatives of $v\mapsto\log(1+v)$
are bounded on $[e^{-M}-1,e^M-1]$. In each positive-order chain-rule
term there is at least one derivative of $e^f-1$. A product of $r$
such derivatives is bounded by
$u^r\le u\max(1,B)^{r-1}$. At order zero, the mean value theorem
between $0$ and $e^{f(y_0)}-1$ gives the same linear bound.
The chain rule is used on the open interval $v>-1$, so no assertion
of a fixed neighborhood contained in that compact interval is needed.

Second, if $\|e^f-1\|_{j,y_0}\le B$ and
$\|h\|_{j,y_0}\le v\le D$, then
\begin{equation}\label{eq:local-parameter-exp-difference}
               \|e^{f+h}-e^f\|_{j,y_0}\le C_{j,B,D}v.
\end{equation}
Use $e^{f+h}-e^f=e^f(e^h-1)$. The zero-order estimate for the second
factor is $|e^{h(y_0)}-1|\le e^D|h(y_0)|$. For positive orders,
the partition formula for derivatives of $e^h$ gives finitely many
products of derivatives of $h$, each containing at least one factor.
Thus $\|e^h-1\|_{j,y_0}\le C_{j,D}v$, by
$v^r\le v\max(1,D)^{r-1}$. Apply
\eqref{eq:local-parameter-product-jet} and
$\|e^f\|_{j,y_0}\le B+1$.

\subsection{The four elementary kinds of factor}
We compare every neighboring factor after the noise rescaling.
For a packet $s$, write its base compensated factor as $w_s$.
The exact ratio identity is
\[
 w_s'-w_s=w_s(e^{H_s}-1),\qquad
 H_s=-(\beta'-\beta)E_s
       -n_s\{\log m(\beta')-\log m(\beta)\}.
\]
Here $E_s$ is its compensated action including the fixed far-field
term, and $n_s$ is the number of chart normalizers in the packet.
The restrictions deciding whether the packet occurs are the same in
both representations. If its enforcement cutoff vanishes on a
neighborhood of $y_0$, both germs vanish. Otherwise the packet action
estimate applies. Let $s_* $ be the maximum packet size and let $m_*$
be the total number of fine bonds. Both are finite on this torus.
Then, with $Y=1+A s_*+m_*$,
\[
 \|H_s\|_{j,y_0}\le tY,\qquad
 \|w_s\|_{j,y_0}\le B.
\]
Choose $t\le Y^{-1}$. The exponential comparison and product rule
give $\|w_s'-w_s\|_{j,y_0}\le C_{\rm p}t$, uniformly over the finite
packet catalogue and all its admissible auxiliary choices.

For a forced incoming factor, the held fine field is the same and
the enforcement cutoff is the same. Its value difference is at most
$b_i$, and its cutoff has a common jet bound $C_{\rm e}$. Hence its
factor difference has jet bound $C_{\rm e}b_i$. This differentiates
only the cutoff, not the incoming forced activity.

For logarithmic factors write $w_i=e^{\ell_i}-1$. The bounded
base logarithm and the summable base factor prices give
$\|\ell_i\|_{j,y_0}\le C u_i$ by
\eqref{eq:local-parameter-log-recovery}. The actual logarithmic
catalogue has two cases. An inherited regular activity has logarithm
minus its held value times the common enforcement cutoff; its
logarithm difference therefore has jet bound $C_{\rm e}a_i$.
Every other catalogue entry is either unchanged by the rescaling,
or is multiplied by $\lambda^2$. This follows term by term from the
quadratic noise terms and $\beta\lambda^2=\beta'$; the path and loop
terms without that quadratic scaling remain unchanged. Thus
\[
 \|\ell_i'-\ell_i\|_{j,y_0}\le C\{\widetilde a_i+t u_i\},
\]
where $\widetilde a_i$ places the regular difference amplitudes at
their inherited logarithmic addresses and is zero elsewhere.
The address embedding is injective, so
$\sum_i\widetilde a_i\le t$. Since $t\le1$ and $u_i\le B$, the
right side is at most $C(1+B)$. Applying
\eqref{eq:local-parameter-exp-difference} gives
\[
                 \|w_i'-w_i\|_{j,y_0}
                     \le C_{\rm l}(\widetilde a_i+t u_i).
\]
Finally, warning factors are exactly unchanged: their charge sets,
forced supports and physical field agree. Their difference price is
zero.

There are finitely many packet labels, say $n_*$. Combining the four
cases gives nonnegative difference prices $v_i$ satisfying
\begin{equation}\label{eq:local-parameter-elementary-differences}
 \|w_i'-w_i\|_{j,y_0}\le v_i,\qquad
 \sum_i v_i\le tK_*,\quad
 K_*=n_*C_{\rm p}+C_{\rm e}+C_{\rm l}(1+B).
\end{equation}
The constants are uniform over the auxiliary choices at this fixed
torus. This is stronger than continuity of each individual factor:
the whole catalogue has a summable difference bound.

\subsection{Products over components}
Set $a_i=u_i+v_i/t$. Then $\sum_i a_i\le B+K_*$,
$\|w_i\|_{j,y_0}\le a_i$, and
$\|w_i'-w_i\|_{j,y_0}\le t a_i$. For every finite set of labels $P$,
put
\[
                    p_P=\prod_{i\in P}3\,3^j a_i.
\]
For $0<t\le1$, induction on $|P|$ yields
\begin{equation}\label{eq:local-parameter-component-products}
 \left\|\prod_{i\in P}w_i'\right\|_{j,y_0}\le p_P,
 \qquad
 \left\|\prod_{i\in P}w_i'-\prod_{i\in P}w_i\right\|_{j,y_0}
                                                     \le t p_P.
\end{equation}
For the induction step use $\|w_i'\|_{j,y_0}\le2a_i$ and
\[
 w_i'G'-w_iG=(w_i'-w_i)G'+w_i(G'-G).
\]
Each product costs at most $3^j$; the chosen factor $3\,3^j a_i$
pays both terms and the bound for the new product. The empty product
has value one and difference zero.

Summing nonnegative products over finite subsets gives
\begin{equation}\label{eq:local-parameter-product-sum}
                 \sum_P p_P\le
                      \exp\{3\,3^j(B+K_*)\}.
\end{equation}
For a finite catalogue this is $\prod_i(1+3\,3^j a_i)$ and
$1+x\le e^x$. Exhaustion and monotone convergence prove the countable
case. Restricting to admissible molecules only decreases the sum.
All supports and selection conditions agree under the common address
identification.

We must identify these products with the actual component integrands.
Assume both representations satisfy the incoming room and activity
bounds for this same sector. The protected compensation map holds
the physical fine field fixed on a neighborhood of $y_0$. On that
neighborhood, the definition of each component therefore equals the
product of the elementary compensated factors with its fixed charge
set. The neighboring compensation has the same protected patch:
its noise course is multiplied by $\lambda$, and its coupling
cancels that factor in the physical field. Applying the same identity
to the neighboring representation proves equality of its component
germ with the primed product. Intersect the two neighborhoods.
Health of the union of the assigned supports implies health of every
factor support, because each is contained in that union. Thus all
the bounds above apply to these actual germs, and their derivatives
are the derivatives of the displayed finite products.

\subsection{The density compensation and the Gaussian ratio}
The exact compensation identity has the anchored form
\[
 K'(y,\lambda z)=K(y,z)
       \exp\{(1-\lambda^2)(q(y,z)-q(y_0,z))\}.
\]
The subtraction at $y_0$ is essential. Put
$h(y)=(1-\lambda^2)(q(y,z)-q(y_0,z))$.
Then $h(y_0)=0$, and every positive-order derivative through $j$
is bounded by $tC F_0(z)^2$. In the partition formula for a derivative
of $e^h$ at $y_0$, a term with $r\ge1$ blocks is bounded by
$t^r C^r F_0(z)^{2r}$. Since $t\le1$ and $F_0\ge1$, it is at most
$tC^r F_0^{2j}$. The zero-order difference $e^{h(y_0)}-1$ is zero.
Consequently, also when $j=0$,
\[
 \|e^h-1\|_{j,y_0}\le tC_jF_0(z)^{2j},\qquad
 \|K'-K\|_{j,y_0}\le tC_1F_0(z)^{N+2j}.
\]
No exponential of an uncontrolled positive polynomial in $z$ is
introduced here.

Write $g_P,g_P'$ for the component products in
\eqref{eq:local-parameter-component-products}. With the fixed
Gaussian density ratio $J_\lambda(z)$, which is independent of $y$,
use the identity
\begin{align*}
 J_\lambda K'g_P'-Kg_P
   ={}&J_\lambda(K'-K)g_P'
       +(J_\lambda-1)Kg_P' +K(g_P'-g_P).
\end{align*}
The preceding estimates and the product rule give
\begin{equation}\label{eq:local-parameter-compensated-bound}
 \|J_\lambda K'g_P'-Kg_P\|_{j,y_0}
 \le t C_2 p_P F_0(z)^{N+2j}
             \{J_\lambda+|J_\lambda-1|/t+1\}.
\end{equation}
Here $C_2$ is fixed and $\sum_PC_2p_P<\infty$ by
\eqref{eq:local-parameter-product-sum}. This is the required local
compensated comparison. The Gaussian estimates then integrate it
with a common integrable dominator and preserve all jets through
order $j$.

\subsection{Packet values outside the healthy jet domain}
For the forced output one also needs value estimates for packets
outside the healthy domain used above. An upper bound on the packet
energy itself is neither available nor needed. The actual packet
has the form
\[
             P_\beta=m_0\exp\{-\beta u+v_\beta\},
 \qquad 0\le m_0\le M_0,\quad u\ge-d_0,
 \quad v_\beta\le A_0,
\]
with $d_0,M_0\ge0$ fixed on this torus. The packet energy lower
bound supplies $d_0$; bounded cutoff and charge factors supply $M_0$.
The unchanged condensation and chart terms, and the finite number
of normalizers, give
$|v_{\beta'}-v_\beta|\le n_0t$ for a finite $n_0\ge0$.
Assume $|\beta'-\beta|\le t\le\min(1,\beta/2)$ and put
\[
 c=\beta/2>0,\qquad
 D=A_0+n_0+(\beta+1)d_0.
\]
If $-d_0\le u\le0$, both exponential arguments are at most $D$,
and their difference is at most $t(d_0+n_0)$. The mean value theorem
gives a bound $t m_0 e^D(d_0+n_0)$.
If $u>0$, both arguments are at most $D-cu$ and their difference
is at most $t(u+n_0)$. Thus
\[
 |P_{\beta'}-P_\beta|
 \le t m_0e^D(u+n_0)e^{-cu}
 \le t m_0e^D(c^{-1}+n_0),
\]
using $cu\le e^{cu}$ and $e^{-cu}\le1$. Combining both cases yields
\begin{equation}\label{eq:local-parameter-forced-packet}
 |P_{\beta'}-P_\beta|
 \le tM_0e^D(d_0+2n_0+c^{-1}).
\end{equation}
Zero packet masks remain zero under the common data and require no
estimate. This proves the packet value comparison uniformly over
the unconstrained positive-energy tail. The corresponding forced
product and Gaussian integration arguments use value bounds; they
do not assert derivative estimates on that tail.

For completeness, the other unrestricted value factors have the same
summable comparison. Let $r_i\ge0$ bound the base logarithms and put
$R_0=\sum_i r_i<\infty$. Their differences are bounded by
$d_i=t r_i+\widetilde a_i$, so $\sum_i d_i\le t(R_0+1)$.
Both the base and new logarithms are at most $2R_0+1$ when $t\le1$.
The mean value theorem gives the exponential-factor difference price
$e^{2R_0+1}d_i$. Forced incoming factors have difference price $b_i$,
since their common cutoff lies in $[0,1]$; warnings cancel. Together
with \eqref{eq:local-parameter-forced-packet}, the finite packet
catalogue therefore has an overall elementary difference budget
$tC_3$.

The corresponding base value prices are summable as well. Packets
and warnings have finite catalogues. The logarithmic catalogue has
its summable base price, while the forced incoming prices are bounded
by their rooted weighted prices. Every forced support is nonempty:
it contains the parent of an inspected incoming inventory site.
Summing the rooted bound over the finitely many torus sites therefore
bounds the unrooted sum of forced prices. The weights removed in this
argument are at least one.

Apply the product argument above with $j=0$ to these value prices.
There are summable $p_P^{\rm v}\ge0$ such that
$|G_P'|\le p_P^{\rm v}$ and
$|G_P'-G_P|\le t p_P^{\rm v}$ for the rescaled full component values.
Before integration use
\[
 |J_\lambda G_P'-G_P|
       \le |J_\lambda-1|p_P^{\rm v}+t p_P^{\rm v}.
\]
The Gaussian ratio estimate at moment order zero bounds the integral
by $tC_4p_P^{\rm v}$. Haar variables keep their normalized measure.
The common support and charge data also identify all eligibility
conditions. A normalized finite tag average preserves this bound;
if a component is ineligible, both terms are zero. This proves the
forced value comparison needed for the exponential expansion of the
output array.
